\section{Inputs, Requirements and Assumptions}
\label{sec:inputs}

\subsection{What is fixed}

Weisshorn is the 2027 EuRoC L9 vehicle (\SI{9}{km} apogee). Inputs come from the LVO, LVS, requirement
database and earlier DJF/DDF (Table~\ref{tab:inputs}).

\begin{table}[H]
\centering
\caption{Inputs from the ERT wiki, the Weisshorn requirements database and the Firehorn documentation.}
\label{tab:inputs}
\small
\begin{tabular}{@{}>{\raggedright\arraybackslash}p{3.2cm}>{\raggedright\arraybackslash}p{4.9cm}>{\raggedright\arraybackslash}p{7.6cm}@{}}
\toprule
Quantity & Value & Source \\
\midrule
External diameter & \SI{180}{mm} & task; 2027\_C\_SE\_REQ\_03 \\
Mission & EuRoC 2027, L9, \SI{9}{km} apogee & Weisshorn wiki page; 2027\_C\_SE\_LVO \\
LV length & $< \SI{5.4}{m}$ (LVO: $\sim\SI{5.4}{m}$) & 2027\_C\_SE\_REQ\_04; 2027\_C\_SE\_LVO \\
Dry / wet mass & $\sim\SI{45}{kg}$ / $\sim\SI{75}{kg}$ ($<\SI{50}{kg}$ dry) & 2027\_C\_SE\_LVO; 2027\_C\_SE\_REQ\_06 \\
Engine & EthaLOx, \SI{4.5}{kN}, \SI{9.1}{s}, $I_t = \SI{40960}{N.s}$, $I_{sp}=\SI{200}{s}$ & 2027\_C\_SE\_LVO; 2027\_C\_SE\_LVS\_DJF \\
Maximum speed & Mach 1.4 & 2027\_C\_SE\_LVO \\
Stability & 1.5--6 cal (EuRoC); FD target 2--4 cal & 2027\_C\_SE\_REQ\_09; 2027\_C\_SE\_FD\_REQ\_07 \\
Launch rail & $84\pm1^\circ$ elevation & 2027\_C\_SE\_FD\_REQ\_05 \\
Payload & 3U CubeSat in the nosecone bay, nose separates at apogee & 2027\_C\_SE\_ST\_FUA; Structure wiki page \\
Nose length (sizing) & scaled from Firehorn: $960\times180/240 = \SI{720}{mm}$ & 2027\_C\_SE\_LVS\_DJF (``scaled down in the same way for mass and length'') \\
Firehorn nose mass & as-built \SI{500}{g} (\SI{960}{mm}, \O\SI{240}{mm}) & 2024\_C\_ST\_NOSECONE\_DJF/DDF; 2025\_C\_ST\_NOSECONE\_PURPOSE\_IR \\
\bottomrule
\end{tabular}
\end{table}

\noindent With $M_{\max} \approx 1.4$ the nose is designed for the whole transonic and low-supersonic
range, and the profile decides how much of the 3U ($100\times100\times\SI{340.5}{mm}$) fits inside.

\subsection{What is assumed}

None of these assumptions changes the ranking.
\begin{itemize}[leftmargin=*,itemsep=2pt]
\item \textbf{Trajectory.} Point mass, rail \SI{12}{m} at $84^\circ$, ISA, no wind, \SI{4.5}{kN} for
\SI{9.1}{s}; propellant $I_t/(I_{sp}g_0) = \SI{20.9}{kg}$ (burn-out \SI{54.1}{kg}). The drag of everything
except the nose is scaled by one factor $k$ so that the baseline nose reaches \SI{9.0}{km}
(Section~\ref{sec:flight}), giving $M_{\max}=1.47$.
\item \textbf{Hot-day envelope.} LVO masses ($75 \rightarrow \SI{45}{kg}$) on a \SI{30}{\celsius} day:
$M_{\max} \approx 1.56$ (a robustness case, Section~\ref{sec:robust}).
\item \textbf{Fins} (not yet defined): four trapezoidal fins, root \SI{300}{mm}, tip \SI{120}{mm}, span
\SI{160}{mm}, sweep \SI{180}{mm}, \SI{5}{mm} thick; they only enter the common drag and the CP shift.
\item \textbf{Payload envelope.} A 3U needs $r \geq \SI{77}{mm}$ (\SI{70.7}{mm} half-diagonal +
\SI{3}{mm} clearance + \SI{2}{mm} guide + \SI{1.2}{mm} shell).
\item \textbf{Nose mass.} Wetted area $\times$ \SI{0.98}{kg/m^2} (Firehorn as-built) + \SI{60}{g} tip.
\item \textbf{Transonic wave drag} ($0.8 < M < 1.1$): smooth rise from zero at $M = 0.8$ to the
$M = 1.1$ value at $M = 1.0$ (Section~\ref{sec:transonic}).
\end{itemize}

\subsection{What could not be found}
\label{sec:flags}

No Weisshorn nosecone requirement exists yet; the Weisshorn fins, CG and mass breakdown are not defined
(so the static margin is given as a change); no supersonic flight data exists at ERT to calibrate drag.
